\Needspace{13\baselineskip}
\chapterrule
\section{Token 与 ID：文字怎样进入数字世界？}
\label{intro:20a-tokenizer}

\phantomsection\label{concept:token}

矩阵接收的是数字，但语言模型首先拿到的是文字。
一句中文如何被拆成离散的 Token，再变成可查表的整数 ID？

\subsection{Token、词表、整数 ID 分别做什么？}
\phantomsection\label{concept:vocabulary}

\textbf{Token} 是文本处理单位，不一定恰好等于自然语言里的词：
它的单位包括单字、子词和标点。
\textbf{Tokenizer（分词器）}按照某套规则把文本拆成 Token，
然后按\textbf{词表}将它们映射为整数 \textbf{ID}。

本书统一使用一份\textbf{人为指定}的五项教学词表：
\[
\begin{array}{c|ccccc}
\text{ID}&1&2&3&4&5\\\hline
\text{Token}&\text{其他}&\text{春天}&\text{。}&\text{花}&\text{开了}
\end{array}
\]
所以“春天 / 花 / 开了 / 。”映射到 $[2,4,5,3]$。
整数编号只标识身份；“开了”和“春天”的重要性要由任务与上下文决定。

\subsection{每个字都拆开，还是把常见片段合起来？}

前面的“春天／花／开了”是人为指定的切分。面对大量文本，怎样自动决定哪些字应该放在一起？
先看“今天天气好”：逐字切开是“今／天／天／气／好”，共 $5$ 个 Token；
如果词表收录“今天”“天气”，同一句话就能写成“今天／天气／好”，共 $3$ 个 Token。
常见片段合在一起，可以让模型少处理几个输入位置。

但若想提前为所有词、词组甚至整句都准备一个 ID，词表就会不断膨胀，
还可能遇到从未收录的新组合。我们希望\textbf{常见组合用整块，其他组合仍能用已有的小片段拼出来}。
一个办法是：保留基础片段，把经常相邻出现的两个片段合成一个新 Token。

先看合并规则怎样使用。假设已经从文本中学到两条规则，按顺序执行：
规则~1：今+天$\to$今天；规则~2：天+气$\to$天气。
切分时先逐字拆开，再按顺序执行：
\[
[\text{今},\text{天},\text{天},\text{气},\text{好}]
\ \xrightarrow{\ \text{规则 1}\ }\
[\text{今天},\text{天},\text{气},\text{好}]
\ \xrightarrow{\ \text{规则 2}\ }\
[\text{今天},\text{天气},\text{好}].
\]
五个片段变成三个 Token。
规则只在文本中真的出现对应相邻对时才生效：
若基础词表也收录了“明”，那么“明天天气好”里没有“今+天”，规则~1 不起作用，
规则~2 仍可使用，结果是 $[\text{明},\text{天},\text{天气},\text{好}]$，共 $4$ 个 Token。

\phantomsection\label{concept:bpe}
这些规则从哪里来？一种方法是在训练文本里统计相邻片段对，
\textbf{每轮把出现次数最多的一对合成新 Token，加入词表，再重新统计}。
这种逐步合并的方法叫作 \textbf{BPE（Byte Pair Encoding，字节对编码）}。
它可以学出单字、词的一部分或整词等不同长度的片段，常被称为\textbf{子词切分}；
这里的“子词”并不表示每个 Token 都恰好是半个词。

这包含两个阶段：\textbf{训练（造词表）}时反复统计、合并，直到词表达到预定大小，
留下词表与有先后顺序的规则清单；\textbf{切分（处理新文本）}时按这份清单合并基础片段，
不再统计频率。下一小节会手算训练时的计数与合并。

现在再看两个数的区别：\textbf{词表大小}是可查表的 Token 条目数，
本例新增了“今天”“天气”两项，基础片段仍保留；
\textbf{序列长度}是这段输入实际切出的 Token 数，“今天天气好”由 $5$ 项变成 $3$ 项。
词表多收录了常见片段，这句话的输入位置反而减少了。

名字里的“字节对”来自它的起源：它原本是一种数据压缩方法，
反复把最常见的一对字节替换成一个新符号。
今天的大模型多采用\textbf{字节级 BPE}，初始片段是 UTF-8 字节而非汉字；
常用汉字在 UTF-8 中占三个字节，一开始就是三个片段。
常见字会被合并回整字甚至词，少见字可能停留在字节片段。
\textbf{因此上面从汉字出发的例子只是教学示意}；
真实的词表和合并规则因模型而异，中文可能被切成字、子词或字节片段。

Token 化影响上下文成本，却\textbf{不会自动赋予 ID 语义}。


\subsection{［复算］五行手算：BPE 怎样挑最频繁的相邻对？}

为了避免中文编码和字节级分词的额外细节，暂把训练语料设成
两段由空格分开的字符序列 \texttt{a b a b} 与 \texttt{a b}，
\textbf{不允许跨训练段落合并}。
\Needspace{14\baselineskip}
每轮统计全部段落内部的\textbf{相邻 Token 对}，合并出现次数最多的一对，
并将新 Token 加入词表。这里没有并列最高频的歧义：

\begin{center}\small
\tagpdfsetup{table/header-rows={1}}
\begin{tabularx}{\linewidth}{@{}p{11mm}p{47mm}X@{}}
\toprule
步骤 & 语料状态或计数 & 词表变化与决策\\\midrule
0 & \texttt{a b a b}；\texttt{a b} & 初始词表 $\{a,b\}$\\
1 & \texttt{(a,b)} 共 3 次；\texttt{(b,a)} 共 1 次 & 选最频繁对 \texttt{a+b}\\
2 & \texttt{ab ab}；\texttt{ab} & 合并为 \texttt{ab}；加入词表\\
3 & \texttt{(ab,ab)} 共 1 次 & 唯一可再合并的相邻对\\
4 & \texttt{abab}；\texttt{ab} & 合并为 \texttt{abab}；加入词表\\
\bottomrule
\end{tabularx}
\end{center}

合并后词表成为 $\{a,b,ab,abab\}$；它\textbf{不会删除初始字符}
（这样未被合并的其他文本仍可编码）。
训练产出的规则清单依次是规则~1：\texttt{a+b}；规则~2：\texttt{ab+ab}。
切分新文本 \texttt{a b a b a b} 时，规则~1 得到 \texttt{ab ab ab}，
规则~2 从左到右合并得到 \texttt{abab ab}，共两个 Token。
真实字节级 BPE 还包含初始字节表示、预分词和特殊 Token 等实现细节。
\noindent{\sffamily\footnotesize\color{blue}\textbf{［复算］}}
执行 \texttt{python3 scripts/bpe\_demo.py} 可以重放两轮计数与合并。

\subsection{真实中文切分：同一句话进入 Qwen2.5}
\phantomsection\label{sec:chinese-tokenization}
\label{experiment:real-chinese-tokenizer}
把“春天花开了。”交给真实模型 Qwen2.5-0.5B 的分词器（不添加特殊 Token），得到：
\begin{center}\small
\tagpdfsetup{table/header-rows={1}}
\begin{tabular}{@{}lrr@{}}\toprule
逐项解码文本 & 真实 ID & 教学词表 ID\\\midrule
春天 & 105303 & 2\\
花 & 99232 & 4\\
开了 & 101467 & 5\\
。 & 1773 & 3\\\bottomrule
\end{tabular}
\end{center}
本句恰好也切成四项，\textbf{真实 ID 与教学词表完全不同}。
去掉句号后，“春天花开了”恰好是三个 Token，与主例的三位置输入对应；主例仍使用教学 ID $[2,4,5]$。
换一句话仍须重新分词，“春天”或“开了”可能被拆成多个 Token。
真实 ID 从 $0$ 编号，查表直接用这个 ID；本书的五项手算词表从 $1$ 数，写 Python 时减 $1$。
想亲自复核，可运行 \texttt{python3 scripts/chinese\_tokenizer\_demo.py --real}，它只下载分词器、不下载模型权重；
固定修订与实测记录见\hyperref[app:tokenizer-record]{附录 G 的中文分词记录}。

\chapterchecktasks{20a-tokenizer}{理解与复算}{
\item \textbf{解释：}为什么整数 Token ID 只是词表索引，而不是“数字越大、词义越强”的语义分数？\chapterinlinehint{想一想重新排列词表编号是否会改变词义。}
\item \textbf{手算：}从五个字符“今／天／天／气／好”出发，假设 BPE 按规则依次合并“今+天”“天+气”。合并后有多少个 Token？列出切分结果。\chapterinlinehint{按给定顺序执行两次合并，并保留未合并项。}
\item \textbf{找错：}同学认为“BPE 一次合并就会自动在所有句子里产生相同的 Token 数”。用词表和输入文本的关系指出漏洞。\chapterinlinehint{固定合并规则，对照不同输入中的相邻片段。}
}


\par\smallskip
Token 化将原始文字变成离散 ID，但 ID 只是类别标签。
下一章先把句子变成预测题：模型看到哪些 ID，每个位置应该预测谁？
随后再计算预测的好坏，并为每个 ID 建立可学习的数值特征。
